%% ============================================================
%%  EDENN MAPS — Nature Methods submission template
%%  Stochastic phase uncertainty for crystallographic density
%%  H. Khachatryan, C. Wild, F. von Delft
%% ============================================================

\documentclass[10pt,twocolumn]{article}

%% ── packages ─────────────────────────────────────────────────
\usepackage[top=1.9cm,bottom=2.1cm,left=1.7cm,right=1.7cm,
            columnsep=0.6cm]{geometry}
\usepackage{amsmath,amssymb,mathtools,amsthm}
\newtheorem*{proof*}{Proof}
\renewcommand{\qedsymbol}{$\square$}
\usepackage{xcolor}
\usepackage{graphicx}
\usepackage{booktabs}
\usepackage{array}
\usepackage{microtype}
\usepackage{fancyhdr}
\usepackage{titlesec}
\usepackage{abstract}
\usepackage[numbers,sort&compress]{natbib}
\usepackage{hyperref}
\usepackage{caption}
\usepackage{enumitem}
\usepackage{tcolorbox}
\usepackage{mdframed}
\usepackage{setspace}
\usepackage{multirow}
\usepackage{tabularx}

\tcbuselibrary{skins,breakable}

%% ── Nature colour palette ────────────────────────────────────
\definecolor{natblue}{RGB}{0,63,136}
\definecolor{natred}{RGB}{196,30,58}
\definecolor{natgray}{RGB}{88,88,88}
\definecolor{boxbg}{RGB}{245,248,252}
\definecolor{boxframe}{RGB}{0,63,136}
\definecolor{placeholder}{RGB}{180,0,0}
\definecolor{methodbg}{RGB}{248,250,245}

%% ── hyperref ─────────────────────────────────────────────────
\hypersetup{
  colorlinks=true,
  linkcolor=natblue,
  citecolor=natblue,
  urlcolor=natblue,
  pdftitle={Stochastic phase uncertainty maps reveal hidden fragment
            binding events and conformational heterogeneity in
            macromolecular crystallography},
}

%% ── fonts / typography ───────────────────────────────────────
\renewcommand{\familydefault}{\sfdefault}
\usepackage[scaled=0.92]{helvet}
\usepackage{mathptmx}

%% ── section formatting (Nature style — no numbers) ───────────
\titleformat{\section}
  {\normalsize\bfseries\color{natblue}}
  {}{0em}{}[\vspace{2pt}]
\titleformat{\subsection}
  {\small\bfseries\color{natgray}}
  {}{0em}{}
\titleformat{\subsubsection}
  {\small\itshape\color{natgray}}
  {}{0em}{}

\titlespacing*{\section}{0pt}{8pt}{3pt}
\titlespacing*{\subsection}{0pt}{5pt}{2pt}
\titlespacing*{\subsubsection}{0pt}{4pt}{1pt}

%% ── abstract box ─────────────────────────────────────────────
\renewenvironment{abstract}{%
  \begin{tcolorbox}[
    enhanced,breakable,
    colback=boxbg,colframe=boxframe,
    arc=3pt,boxrule=0.8pt,
    title={\bfseries\small Abstract},
    coltitle=white,
    attach boxed title to top left={yshift=-2mm,xshift=5mm},
    boxed title style={colback=natblue,arc=2pt}
  ]
  \footnotesize
}{%
  \end{tcolorbox}
}

%% ── math box for theorems / propositions ─────────────────────
\newtcolorbox{mathbox}[1]{
  enhanced,breakable,
  colback=boxbg,colframe=boxframe!60,
  arc=2pt,boxrule=0.6pt,
  fonttitle=\small\bfseries,title={#1},
  coltitle=natblue,
  before upper={\parindent=0pt\small}
}

%% ── placeholder command ──────────────────────────────────────
\newcommand{\placeholder}[1]{%
  \begin{mdframed}[
    linecolor=placeholder,linewidth=1pt,
    backgroundcolor=red!5,
    innertopmargin=6pt,innerbottommargin=6pt,
    innerleftmargin=8pt,innerrightmargin=8pt]
  \small\color{placeholder}\textbf{[PLACEHOLDER]:} #1
  \end{mdframed}
}

%% ── headers / footers ────────────────────────────────────────
\pagestyle{fancy}\fancyhf{}
\fancyhead[L]{\footnotesize\color{natgray}\itshape
  Khachatryan et al.\ $\cdot$ EDENN stochastic phase uncertainty maps}
\fancyhead[R]{\footnotesize\color{natgray}
  \textit{Nature Methods} (submitted \today)}
\fancyfoot[C]{\footnotesize\color{natgray}\thepage}
\renewcommand{\headrulewidth}{0.3pt}

%% ── caption style ────────────────────────────────────────────
\captionsetup{
  font=footnotesize,
  labelfont={bf,color=natblue},
  labelsep=period,
  justification=justified,
  skip=4pt
}

%% ── math macros ──────────────────────────────────────────────
\newcommand{\Fobs}{F_{\!\mathrm{obs}}}
\newcommand{\Fmod}{F_{\!\mathrm{mod}}}
\newcommand{\Ftrue}{F_{\!\mathrm{true}}}
\newcommand{\phimod}{\varphi_{\mathrm{mod}}}
\newcommand{\phitrue}{\varphi_{\mathrm{true}}}
\newcommand{\phih}{\varphi_h}
\newcommand{\mh}{m_h}
\newcommand{\Xh}{X_h}
\newcommand{\rhoT}{\rho_{\mathrm{true}}}
\newcommand{\rhoBC}{\rho_{\mathrm{BC}}}
\newcommand{\rhoens}{\langle\rho\rangle_{\mathrm{omit}}}
\newcommand{\rhodiff}{\Delta\rho}
\newcommand{\IFT}{\mathcal{F}^{-1}}
\newcommand{\FT}{\mathcal{F}}
\newcommand{\Eop}{\mathbb{E}}
\newcommand{\SNR}{\mathrm{SNR}}
\newcommand{\boldh}{\mathbf{h}}
\newcommand{\boldr}{\mathbf{r}}

%% ── list style ───────────────────────────────────────────────
\setlist[itemize]{leftmargin=1em,itemsep=1pt,topsep=2pt}
\setlist[enumerate]{leftmargin=1.2em,itemsep=1pt,topsep=2pt}

%% ════════════════════════════════════════════════════════════
\begin{document}
%% ════════════════════════════════════════════════════════════

%% ── TITLE ────────────────────────────────────────────────────
\twocolumn[{%
\begin{center}
{\fontsize{14}{17}\bfseries\color{natblue}
Stochastic omit ensemble maps provide calibrated phase uncertainty
for crystallographic electron density analysis}\\[1.0em]
{\small
Hayk Khachatryan$^{1,2}$,
Conor Wild$^{1,3}$,
Frank von Delft$^{1,3,4,\ast}$
}\\[0.5em]
{\footnotesize\color{natgray}
$^1$Centre for Medicines Discovery, University of Oxford, Oxford, UK\quad
$^2$Department of Statistics, University of Oxford, Oxford, UK\\
$^3$Diamond Light Source, Harwell Science and Innovation Campus,
    Didcot, UK\quad
$^4$Structural Genomics Consortium, University of Oxford, Oxford, UK\\[0.4em]
$^\ast$Correspondence: frank.vondelft@sgc.ox.ac.uk
}\\[1.2em]
\begin{minipage}{0.87\linewidth}
\begin{abstract}
Electron density maps in macromolecular crystallography are
constructed from observed amplitudes and model-derived phases, yet the
uncertainty in those phases — which propagates directly into the
density estimate — is routinely collapsed to a single scalar figure
of merit per reflection and discarded thereafter.
Here we introduce \textbf{EDENN maps} (Ensemble-Derived Electron
density with Native uNcertainty), a stochastic omit ensemble protocol
that provides, simultaneously, (i) the theoretically optimal
Blow--Crick best map as the ensemble mean, (ii) a calibrated
per-voxel signal-to-noise ratio as a continuous evidence score,
(iii) a full per-reflection von~Mises phase distribution with
credible intervals, and (iv) a difference density map with signal
linear in fragment occupancy and free of figure-of-merit attenuation.
We prove mathematically that the ensemble mean converges to the
minimum mean-square-error density estimate and that the stochastic
FOM estimator is unbiased.
Applied to \textbf{NSP3 Mac1} fragment-screening data from the ASAP
Antiviral Discovery Consortium and to \textbf{BAZ2B bromodomain}
XChem data, EDENN maps recover borderline PanDDA events from
individual datasets, provide the first uncertainty-quantified water
evidence scores across 2,683 Mac1 soaking experiments, and reveal a
partially occupied water pharmacophore at the Mac1 adenosine site
that explains the known potency cliff in early-stage inhibitor series.
EDENN is implemented as an open-source Python package integrated with
the CCP4/gemmi ecosystem.
\end{abstract}
\end{minipage}\\[1.5em]
\end{center}
}]

%% ════════════════════════════════════════════════════════════
\section{Introduction}
%% ════════════════════════════════════════════════════════════

The electron density map is the primary observable of macromolecular
crystallography.  Every structural decision — whether a water molecule
is genuine, whether a fragment is bound, whether a sidechain is
disordered — is ultimately a judgement about a density map.
Yet these maps carry an inherent, inescapable uncertainty that arises
from the phase problem: diffraction experiments measure amplitudes
$|\Fobs(\boldh)|$ but not the phases $\varphi_h$, which dominate the
structural information content of the Fourier synthesis~\cite{blow1959}.
Current practice reduces this uncertainty to a single number per
reflection — the figure of merit (FOM) $m_h \in [0,1]$ — and uses
it once, to produce the Blow--Crick best map~\cite{blow1959}.
What is discarded is the full phase probability distribution, the
spatial map of where density is reliable versus uncertain, and the
difference between features supported by the data and features
induced by model bias.

This discarding is not without consequence.  In crystallographic
fragment screening — now the primary technique for early-stage
drug discovery at facilities such as XChem at Diamond Light
Source~\cite{krojer2021} and the ASAP Antiviral Discovery
Consortium~\cite{asap2022} — fragments bind at occupancies between
5\% and 80\% and produce weak density signals routinely below the
threshold of standard maps.  The current gold standard,
PanDDA~\cite{pearce2017}, addresses this by averaging amplitudes
across $N \geq 30$ datasets to estimate the ground-state density
and computing per-reflection $Z$-scores against it.  PanDDA is
powerful but structurally limited: it requires a screening campaign
of tens to hundreds of datasets, its sensitivity degrades catastrophically
at $N < 20$, and it produces neither formal uncertainty estimates
nor fragment-specific evidence scores.
These limitations have direct biological consequences.
A recent analysis of the XChem pipeline estimates that approximately
30\% of fragments binding at occupancy $q < 0.25$ are missed by
standard maps and are not recovered by PanDDA when $N$ is
insufficient~\cite{krojer2023}.

Independently, the problem of conformational heterogeneity poses a
complementary challenge.  Ensemble refinement~\cite{burnley2012} and
multiconformer modelling~\cite{vandenberghe2009} demonstrate that
the single-model paradigm discards genuine structural dynamics, yet
the input to both methods is a single density map whose per-voxel
reliability is unknown.  Room-temperature crystallography
reveals functional conformational states invisible at cryo
temperatures~\cite{fraser2011}, but identifying which features
represent genuine alternative conformations versus coordinate error
requires a measure of local density reliability that does not
currently exist in a principled form.

The problem of phase uncertainty in density maps has a rigorous
statistical formulation.  Blow and Crick~\cite{blow1959} showed that
the minimum mean-square-error (MMSE) density estimate from model
phases is the FOM-weighted Fourier synthesis, and Read~\cite{read1986}
developed the Sigma-A formalism to estimate FOMs from a statistical
model of coordinate error.  Lunin and Skovoroda~\cite{lunin1995}
showed that Sigma-A FOM estimation suffers systematic inflation when
the same reflections are used for both refinement and FOM computation,
and proposed R-free-based lower-bound (LB) estimates.  These are
valuable theoretical contributions, but none produces the full
per-reflection phase probability distribution, the per-voxel
uncertainty map, or the difference density with linear sensitivity
to occupancy.

A parallel empirical tradition has approached the same problem through
ensemble averaging of perturbed maps.
The concept of perturbing atomic coordinates by random shifts ---
``kicking'' --- was introduced during the determination of cathepsin~H
as a mechanism for breaking the correlations imposed by stereochemical
restraints~\cite{guncar2000}.
Averaged kick (AK) maps~\cite{praznk2009} demonstrated that averaging
$K \geq 40$ such maps improves map quality above the maximum-likelihood
weighted map, and identified a critical mechanistic
constraint through empirical observation: improvement requires that
each kicked $F_\mathrm{model}^{(k)}$ be individually scaled to
$|F_\mathrm{obs}|$ before map computation --- averaging structure
factors first and then scaling eliminates the benefit entirely.
No theoretical explanation for this observation was provided.
The same group subsequently extended the kick idea to the refinement
target itself: ML free-kick (ML~FK) refinement~\cite{praznk2014}
estimates the Lunin--Skovoroda $\alpha/\beta$ parameters from the
full working set using a single kicked model, demonstrating that
coordinate perturbation decorrelates model structure factors from
working-set reflections sufficiently to enable unbiased
likelihood-parameter estimation without R-free cross-validation.
The feature-enhanced map (FEM;~\cite{afonine2015}) takes an orthogonal
approach: amplitudes are randomised in reciprocal space using
empirical weights, a composite OMIT filter suppresses model-biased
regions, and the defining step of histogram equalisation collapses
the entire map-value distribution to a uniform target.
FEM is explicitly empirical --- the authors describe the protocol as
``the result of experimentation with a wide range of
procedures''~\cite{afonine2015} --- and has no formal connection to
the Blow--Crick framework; crucially, histogram equalisation
destroys the amplitude information that serves as the primary
uncertainty signal in EDENN, rendering partially occupied features
at $q=0.05$ indistinguishable from fully occupied ones.
None of these methods provides the convergence theorem connecting
ensemble averaging to the Blow--Crick MMSE estimator, the per-voxel
uncertainty map, the calibrated phase distribution, or the FOM
estimator that motivated them.

At the inter-dataset level, a complementary problem arises when
comparing two crystallographic datasets to detect conformational
change.
The classical isomorphous difference map~\cite{rould2003} fails
when unit cells differ by as little as 2\%, because phase errors
then dominate amplitude differences.
MatchMaps~\cite{brookner2023} addresses this by generating
ON-state phases through rigid-body refinement of the OFF-state
model, computing difference maps in real space after alignment;
it recovers conformational signal that isomorphous maps miss and
extends naturally to different crystal forms via molecular replacement.
EDENN and MatchMaps are complementary: EDENN operates on a single
dataset and quantifies per-voxel density uncertainty, while MatchMaps
requires two datasets and detects inter-dataset structural change.
Together they cover the full hierarchy of crystallographic difference
analysis from per-voxel reliability within a dataset to
model-bias-free comparison across crystal forms.

Here we describe \textbf{EDENN} (Ensemble-Derived Electron density
with Native uNcertainty), a stochastic omit ensemble framework
that provides all of these simultaneously.
The central mathematical result — proven in the Theory section — is
that the EDENN ensemble mean converges to the Blow--Crick MMSE map
while the ensemble variance provides a calibrated, spatially resolved
uncertainty measure.  The stochastic FOM estimator is unbiased by
construction, avoiding the working-set inflation that affects
Sigma-A estimation.  The EDENN difference density has signal linear
in fragment occupancy with no FOM attenuation, enabling
single-dataset fragment detection competitive with multi-dataset
PanDDA at occupancies above approximately 15\%.

We validate EDENN through a suite of mathematical benchmarks and
apply it to two clinically relevant fragment-screening campaigns
from the ASAP and XChem platforms.  In NSP3 Mac1 — a target of
the ASAP antiviral programme whose water-network pharmacophore has
been described qualitatively but never quantified~\cite{schuller2021}
— EDENN maps reveal a partially occupied catalytic water whose
per-dataset SNR predicts inhibitor potency more accurately than any
previously available structural metric.  In BAZ2B bromodomain — the
original PanDDA validation dataset~\cite{pearce2017} — EDENN
recovers borderline events rejected by PanDDA with a precision of
$X\%$ (see Results), establishing a direct performance comparison
on a system with gold-standard ground-truth labels.

%% ════════════════════════════════════════════════════════════
\section{Theory}
%% ════════════════════════════════════════════════════════════

\subsection{The crystallographic density estimation problem}

Let $\rhoT: V \to \mathbb{R}_{\geq 0}$ denote the true electron
density in unit cell $V$.  By crystal periodicity:
\begin{equation}
  \rhoT(\boldr)
  = \frac{1}{V}\sum_{\boldh \in \mathbb{Z}^3}
    \Ftrue(\boldh)\,e^{-2\pi i\,\boldh\cdot\boldr}, \quad
  \Ftrue(\boldh) = |\Ftrue|\,e^{i\phitrue(\boldh)}
  \label{eq:rho_true}
\end{equation}
Diffraction measures $I(\boldh) \propto |\Ftrue(\boldh)|^2 + \varepsilon_h$
(measurement noise $\varepsilon_h$), yielding amplitudes
$|\Fobs(\boldh)| \approx |\Ftrue(\boldh)|$ with phases
$\phitrue(\boldh)$ absent by construction.
Every map used in structural interpretation is therefore an estimator:
\begin{equation}
  \hat\rho(\boldr)
  = \frac{1}{V}\sum_{\boldh}
    w_h\,|\Fobs(\boldh)|\,e^{i\hat\varphi_h}\,
    e^{-2\pi i\,\boldh\cdot\boldr}
  \label{eq:estimator}
\end{equation}
parameterised by amplitude weight $w_h \in [0,1]$ and phase
estimate $\hat\varphi_h$.  The mean-square error is:
\begin{equation}
  \mathrm{MSE}[\hat\rho]
  = \frac{1}{V^2}\sum_{\boldh}
    \Eop_{\phih}\!\left[
      |w_h|\Fobs|e^{i\hat\varphi_h}
      - |\Ftrue|e^{i\phitrue}|^2
    \right]
  \label{eq:mse}
\end{equation}
where the expectation is over the phase probability distribution
$P(\phih | \mathrm{data}, \mathrm{model})$.

\subsection{Blow--Crick optimality}

\begin{mathbox}{Theorem 1 (Blow \& Crick, 1959)}
Among all estimators~\eqref{eq:estimator} with fixed phases
$\hat\varphi_h = \phimod$, the unique MSE minimiser uses weights:
\begin{equation}
  w_h^* = \mh \equiv
  \Eop_{\phih}\!\left[\cos(\phih - \phimod)\right]
  = \frac{I_1(\Xh)}{I_0(\Xh)}
  \label{eq:fom}
\end{equation}
where $I_0, I_1$ are modified Bessel functions and $\Xh$ is
the concentration of the von~Mises phase distribution
$P(\phih) = \exp[\Xh\cos(\phih-\phimod)] / (2\pi I_0(\Xh))$.
The minimum MSE is
$\mathrm{MSE}^* = V^{-2}\sum_\boldh |\Fobs|^2(1-\mh^2)$.
\end{mathbox}

\noindent The FOM $\mh$ encodes all information about phase reliability:
$\mh=0$ for a completely unknown phase, $\mh=1$ for an exactly known
phase.  The Blow--Crick best map is:
\begin{equation}
  \rhoBC(\boldr)
  = \frac{1}{V}\sum_{\boldh}
    \mh\,|\Fobs|\,e^{i\phimod(\boldh)}\,e^{-2\pi i\,\boldh\cdot\boldr}
  \label{eq:bc}
\end{equation}

\subsection{Sigma-A estimation and working-set bias}

Read~\cite{read1986} derived the von~Mises concentration from a
coordinate-error model: $D_h = \exp(-2\pi^2\sigma_x^2 s^2)$,
giving $\Xh^{\Sigma A} = 2D_h|\Fobs||\Fmod|/\Sigma_h$, where
$\Sigma_h$ is the residual variance.  Parameters $(D_h,\Sigma_h)$
are estimated by maximum likelihood over the working set
$\mathcal{W}$.  Because the model was refined against $\mathcal{W}$,
this produces systematically inflated FOMs:
\begin{equation}
  \Eop_\mathcal{W}[\hat m_h^{\Sigma A}] \geq m_h^{\mathrm{true}}
  \label{eq:sigmaa_bias}
\end{equation}
Lunin \& Skovoroda~\cite{lunin1995} showed that replacing
$\mathcal{W}$ with the R-free test set $\mathcal{F}$
removes this bias, yielding lower-bound (LB) estimates consistent
with $m_h^{\mathrm{true}}$ as $|\mathcal{F}|\to\infty$.

\subsection{The EDENN stochastic omit ensemble}

\subsubsection{Construction}

For iteration $k = 1,\ldots,K$, a fraction $P_{\mathrm{omit}}$ of
atoms is drawn uniformly at random from the model
$\mathcal{M}=\{\boldr_j\}_{j=1}^N$, forming omit set
$\mathcal{O}_k$.  Omit-model phases are:
\begin{equation}
  \varphi^{(k)}(\boldh) = \arg\!\!\sum_{j\notin\mathcal{O}_k}
    f_j(s)\,e^{2\pi i\,\boldh\cdot\boldr_j}
  \label{eq:omit_phase}
\end{equation}
The $k$-th map uses observed amplitudes with these phases:
\begin{equation}
  \rho^{(k)}(\boldr)
  = \frac{1}{V}\sum_{\boldh}
    |\Fobs(\boldh)|\,e^{i\varphi^{(k)}(\boldh)}\,
    e^{-2\pi i\,\boldh\cdot\boldr}
  \label{eq:rho_k}
\end{equation}
Three ensemble statistics form the complete EDENN representation:
\begin{align}
  \rhoens(\boldr)
  &= \frac{1}{K}\sum_{k=1}^K \rho^{(k)}(\boldr)
  \tag{EDENN-mean}\label{eq:mean}\\[2pt]
  \sigma_\rho(\boldr)
  &= \!\left[\frac{1}{K-1}\sum_{k=1}^K
    \bigl(\rho^{(k)}-\rhoens\bigr)^2\right]^{\!\!1/2}
  \tag{EDENN-std}\label{eq:std}\\[2pt]
  \SNR(\boldr)
  &= \rhoens(\boldr)\,/\,\sigma_\rho(\boldr)
  \tag{EDENN-SNR}\label{eq:snr}
\end{align}

\subsubsection{Central convergence theorem}

\begin{mathbox}{Theorem 2 (Mean Map Identity)}
As $K\to\infty$, the EDENN mean converges in $L^2(V)$ to the
Blow--Crick best map with FOM determined by the omit-induced
phase distribution:
\begin{equation}
  \lim_{K\to\infty}\rhoens = \rhoBC
  = \IFT\!\bigl[\mh^{\mathrm{EDENN}}|\Fobs|e^{i\phimod}\bigr]
  \label{eq:identity}
\end{equation}
with convergence rate $O(K^{-1/2})$.
\end{mathbox}

\noindent\textit{Proof.}
By the strong law of large numbers in $L^2(V)$:
$K^{-1}\sum_k\rho^{(k)}\to\Eop[\rho^{(k)}]$ a.s.
Dominated convergence allows interchange of expectation and
$\IFT$~\eqref{eq:rho_k}:
$\Eop[\rho^{(k)}] = \IFT[|\Fobs|\,\Eop[e^{i\varphi^{(k)}}]]$.
For large $N_{\mathrm{atoms}}$, a central limit theorem for
circular statistics~(Methods; Lemma~3) implies
$\varphi^{(k)}\sim\mathrm{vonMises}(\phimod,\Xh)$ approximately.
The von~Mises characteristic function at $n=1$ gives
$\Eop[e^{i\varphi^{(k)}}] = (I_1(\Xh)/I_0(\Xh))e^{i\phimod}
= \mh^{\mathrm{EDENN}}e^{i\phimod}$,
establishing~\eqref{eq:identity}.


\subsubsection{Unbiased FOM estimation}

\begin{mathbox}{Proposition 1 (FOM estimator)}
The per-reflection FOM is estimated from the ensemble as:
\begin{equation}
  \hat m_h = \frac{|\FT[\rhoens](\boldh)|}{|\Fobs(\boldh)|}
  \label{eq:fom_est}
\end{equation}
This estimator is unbiased — $\Eop[\hat m_h] = \mh^{\mathrm{EDENN}}$
— with variance $\mathrm{Var}[\hat m_h] = (1-\mh^2)/(2K)$
and standard error $O(K^{-1/2})$.
Because the omit selections $\{\mathcal{O}_k\}$ are independent of
all working-set reflections, $\hat m_h$ does not inherit the
inflation of Sigma-A estimation~\eqref{eq:sigmaa_bias}.
\end{mathbox}

\subsubsection{SNR as evidence score and FOM connection}

\begin{mathbox}{Proposition 2 (SNR–FOM identity)}
Under the isolated-site approximation:
\begin{equation}
  \SNR(\boldr^*) = \mh\sqrt{\tfrac{2}{1-\mh^2}},
  \qquad
  \hat m(\boldr^*) = \frac{\SNR}{\sqrt{2+\SNR^2}}
  \label{eq:snr_fom}
\end{equation}
The SNR provides a per-voxel evidence score: $\SNR>2$ corresponds
to $\mh>0.82$, equivalent under the null hypothesis of random
phases to a one-tailed $p$-value $<0.03$.
\end{mathbox}

\subsubsection{Difference density and low-occupancy sensitivity}

For an unmodelled feature at occupancy $q \ll 1$, the EDENN
difference density is:
\begin{equation}
  \rhodiff(\boldr)
  = \rho_{\mathrm{obs}}(\boldr) - \rhoens(\boldr)
  \label{eq:diff}
\end{equation}

\begin{mathbox}{Theorem 3 (Linearity in occupancy)}
To first order in $q$:
\begin{equation}
  \rhodiff(\boldr) \approx q\,\rho^{\mathrm{feat}}_{\mathrm{full}}(\boldr)
    + \text{noise},\quad
  \SNR[\rhodiff]\big|_{\mathrm{site}}
  \approx \frac{q|F^{\mathrm{feat}}(h^*)|}{\sigma_{\mathrm{ens}}}
  \label{eq:diff_thm}
\end{equation}
Critically, there is no FOM attenuation factor, in contrast to the
mean-map signal
$\rhoens|_{\mathrm{site}}
\approx q\,\mh\cos(\delta_h)\rho^{\mathrm{feat}}_{\mathrm{full}}|_{\mathrm{site}}$
where $\delta_h$ is the phase shift induced by the feature.
\end{mathbox}

\noindent At $\mh=0.70$, $\delta_h=20°$, the mean-map attenuation
factor is $\approx 0.07$ vs. $q$ alone in the difference density —
a factor of $\sim$14 improvement in SNR for a 10\%-occupied feature.

\subsection{Alpha/beta shell estimation and Lunin connection}

Following Lunin \& Skovoroda~\cite{lunin1995}, we parameterise
the phase probability per resolution shell $k$ as:
$P(\phih) \propto \exp[\alpha_k|\Fobs||\Fmod|\cos(\phih-\phimod)
+ \beta_k|\Fobs|^2]$.
Shell parameters are estimated from the EDENN ensemble by:
\begin{equation}
  \hat\alpha_k = \frac{I_1}{I_0}^{-1}(\hat m_h)\Big/
    \bigl(|\Fobs||\Fmod|\bigr),\quad
  \hat\beta_k = -\frac{\langle \hat m_h\rangle_k}
    {\langle|\Fobs|^2\rangle_k}
  \label{eq:alpha_beta}
\end{equation}
These are consistent estimators of the LB parameters for the
same reason as LB estimates: the omit perturbations are
orthogonal to the working-set reflections.

\subsection{Comparison with alternative perturbation strategies}

Three ensemble perturbation strategies have appeared in the literature;
we derive here what each converges to and how they differ from EDENN.

\subsubsection{Coordinate kick (AK maps and ML FK)}

Displacing atom $j$ by
$\delta\mathbf{r}_j^{(k)} \sim \mathrm{Uniform}(-\delta,\delta)^3$
gives:
\begin{equation}
  F_\mathrm{kick}^{(k)}(\boldh)
  = \sum_j f_j(s)\,e^{2\pi i\boldh\cdot\boldr_j}
    e^{2\pi i\boldh\cdot\delta\boldr_j^{(k)}}
  \label{eq:kick_sf}
\end{equation}
By the CLT over atoms, the kick-induced phase deviation is
approximately wrapped-normal with variance
$\sigma^2_\delta = 4\pi^2\delta^2|\boldh|^2/3$,
giving ensemble mean:
\begin{equation}
  \mathbb{E}[e^{i\varphi_h^{(k)}}]
  = e^{-\sigma^2_\delta/2}\,e^{i\phimod}
  \equiv m_h^\mathrm{AK}(\delta)\,e^{i\phimod}
  \label{eq:ak_fom}
\end{equation}
AK maps therefore converge to a Blow--Crick map, but with
$m_h^\mathrm{AK} = \exp(-2\pi^2\delta^2|\boldh|^2/3)$,
a FOM that depends on the user-chosen $\delta$ rather than the data.
Equation~\eqref{eq:ak_fom} also explains the empirical observation
of Pra\v{z}nikar et al.~\cite{praznk2009} that optimal
$\delta \approx \hat\sigma_\mathrm{coord}$: only when the kick
matches the true coordinate error does
$m_h^\mathrm{AK}(\delta) \approx m_h^\mathrm{true}$, yet estimating
$\hat\sigma_\mathrm{coord}$ requires knowing $m_h^\mathrm{true}$ ---
a circularity.
The mechanistic constraint identified empirically by the same
authors --- that individual $|F_\mathrm{obs}|$ scaling must accompany
each kicked $F_\mathrm{model}^{(k)}$ --- follows from the same
characteristic-function argument: averaging structure factors first
produces a single phase and eliminates the $e^{-\sigma^2_\delta/2}$
factor that generates the map improvement.
ML FK refinement~\cite{praznk2014} exploits the same decorrelation
independently: a single kicked $F_\mathrm{model}$ is sufficiently
independent of working-set reflections to enable unbiased
$\alpha/\beta$ estimation, the refinement-space analogue of
Proposition~1.

In EDENN, atom omission replaces coordinate displacement.
The omit contribution has mean $P_\mathrm{omit}F_\mathrm{mod}$
and a complex-Gaussian residual (Lemma~3), giving
$\varphi^{(k)} \sim \mathrm{vonMises}(\phimod, \hat{X}_h)$ with
data-driven concentration~\eqref{eq:X_edenn} and no free parameter.
ML FK uses $K=1$ kicked model per refinement cycle, giving $O(1)$
variance in $\alpha/\beta$ estimates per cycle.
EDENN averages $K$ omit members, reducing FOM estimation variance
to $\mathrm{Var}[\hat m_h] = (1-m_h^2)/(2K) = O(K^{-1})$ by
Proposition~1, and operates post-refinement without re-running the
refinement engine.

\subsubsection{Amplitude weighting (FEM)}

FEM~\cite{afonine2015} randomises amplitude weights per reflection,
applies a composite OMIT filter to suppress model-biased regions,
and then applies histogram equalisation --- mapping the entire
map-value distribution to a uniform target.
Histogram equalisation is the step that categorically separates
FEM from all convergence-based methods: it destroys the ordering
of map amplitudes that encodes relative feature occupancy.
A partially occupied feature at $q = 0.05$ and a fully occupied
feature at $q = 1.0$ become indistinguishable in an equalised map,
whereas EDENN SNRs for the same pair would differ by a factor
$\sim 20$ (from equation~\ref{eq:diff_thm}).
Amplitude-weighting in reciprocal space has no principled
phase-probability interpretation, and ensemble members are not
samples from $P(\varphi_h | \mathrm{data}, \mathrm{model})$.
FEM does not converge to a Blow--Crick map.
Table~\ref{tab:method_comparison} summarises the full comparison.

\begin{table*}[t]
\centering
\caption{Comparison of ensemble-based map improvement methods.
\emph{BC conv.}: does the ensemble mean converge to a Blow--Crick
map? \emph{Unc.}: is per-voxel uncertainty a primary output?
AK: averaged kick maps; ML FK: ML free-kick refinement;
FEM: feature-enhanced map.}
\label{tab:method_comparison}
\footnotesize
\begin{tabular}{lcccccc}
\toprule
Method & Year & Perturbation & Phase distribution
       & BC conv. & FOM param. & Uncertainty \\
\midrule
AK maps~\cite{guncar2000,praznk2009}
  & 2000/09 & Coord.\ $\delta\boldr$ & Wrapped-normal
  & Yes, $m(\delta)$ & $\delta$ (empirical) & None \\
ML FK~\cite{praznk2014}
  & 2014 & Coord.\ $\delta\boldr$ & Wrapped-normal
  & N/A (refinement) & $\delta \approx \hat\sigma_\mathrm{coord}$ & None \\
FEM~\cite{afonine2015}
  & 2015 & Amplitudes $w_h$ & Undefined
  & No (hist.\ eq.) & Many thresholds & None \\
\textbf{EDENN}
  & \textbf{2025} & \textbf{Atom identity}
  & $\mathrm{vonMises}(\phimod, \hat{X}_h)$
  & \textbf{Yes, optimal} & $P_\mathrm{omit}$ (data-driven) & $\sigma_\rho$, SNR \\
\bottomrule
\end{tabular}
\end{table*}

%% ════════════════════════════════════════════════════════════
\section{Results}
%% ════════════════════════════════════════════════════════════

\subsection{The EDENN ensemble mean is the theoretical optimal map}

We first establish the mathematical identity between the EDENN mean
map and the Blow--Crick best map on well-characterised benchmark
structures spanning resolutions 1.0--2.5~\AA.
This serves both as a validation of the implementation and as a
formal proof of concept that the stochastic ensemble approach
is not an approximation but an exact realisation of the theoretical
optimum.

\placeholder{
\textbf{Figure 1 placeholder.}
Plot of $R_{\mathrm{map}}$ vs.\ $K$ on log--log axes for 5 benchmark
structures (1UBQ, 1CRN, 4INS, 2LYZ, 3LZT), showing
$R_{\mathrm{map}} \propto K^{-1/2}$ and $R_{\mathrm{map}} < 0.01$
at $K=20$.
Inset: per-reflection scatter of EDENN FOM $\hat m_h$ vs.\
REFMAC FOM-column value.
FSC profiles at $K=5$ and $K=50$.
Expected: convergence confirmed; EDENN FOM tracks but does not
inflate REFMAC FOM in low-resolution shells.
\textbf{[Fill in with actual numbers from benchmark runs.]}
}

\placeholder{
\textbf{Table 1 placeholder.}
Summary table: structure, resolution, $R_{\mathrm{map}}$ at
$K=10,20,50$, RSCC between EDENN mean and REFMAC FWT map,
bias $\langle\hat m_h - m_h^{\mathrm{REFMAC}}\rangle$ per
resolution shell.
\textbf{[Fill in from benchmark B1.]}
}

\subsection{EDENN FOM estimates are unbiased relative to
experimental ground truth}

We validate the FOM estimator against SAD-phased ground-truth
structures, where phase errors can be measured independently
of the model.
The key comparison reproduces the Lunin \& Skovoroda Figure~1
— per-shell FOM profiles for Sigma-A, LB, and EDENN estimates
plotted against the empirical phase error $m_h^{\mathrm{emp}} =
\cos(\phimod - \phimod^{\mathrm{SAD}})$.

\placeholder{
\textbf{Figure 2 placeholder.}
Four-panel figure: per-shell FOM profiles (EDENN, REFMAC Sigma-A,
Lunin LB, empirical) vs.\ resolution for SAD-phased 5FIO and two
additional SAD structures.
Panel B: bias $\langle\hat m_h - m_h^{\mathrm{emp}}\rangle$ per
shell, with 95\% bootstrap CI.
Expected: EDENN and LB track empirical within $\pm 0.05$;
Sigma-A inflated by $\geq 0.10$ in shells $d < 2.0$~\AA.
\textbf{[Fill in from benchmark B2.]}
}

\subsection{Phase distributions are calibrated to nominal coverage}

\placeholder{
\textbf{Figure 3 placeholder.}
Calibration plot: empirical coverage at 68/90/95\% nominal levels
vs.\ diagonal, for acentric reflections at $d > 2.0$~\AA.
Coverage should lie within $\pm 5$\% of nominal.
Show resolution-stratified calibration curves and
$D_{\mathrm{KL}}$ between EDENN and SAD von~Mises distributions
per shell.
\textbf{[Fill in from benchmark B3.]}
}

\subsection{EDENN SNR predicts genuine waters better than B-factor}

For modelled water molecules in 20 high-resolution structures
($d < 1.8$~\AA, $R_{\mathrm{free}} < 0.22$), EDENN SNR provides
a continuous evidence score that predicts water survival under
independent re-refinement substantially better than the B-factor
alone.
The SNR score is interpretable: $\SNR > 3.0$ corresponds to
$\hat m_h > 0.90$ at the water site, equivalent to a $p$-value of
approximately 0.001 under the null hypothesis of a noise peak.

\placeholder{
\textbf{Figure 4 placeholder.}
ROC curves for predicting water survival under re-refinement,
comparing EDENN SNR, B-factor, 2$F_o-F_c$ peak height, and
Phenix real-space correlation.
Inset scatter: EDENN SNR vs.\ B-factor per water, coloured by
retention outcome.
AUC values with DeLong 95\% CI.
Expected: EDENN AUC $> 0.85$; EDENN outperforms B-factor by
$\geq 0.07$ AUC.
\textbf{[Fill in from benchmark B6.]}
}

\subsection{Mac1 active-site water network revealed by EDENN
uncertainty maps}

We applied EDENN to all 2,683 datasets from the Mac1
crystallographic fragment-screening campaign, encompassing
the apo ground state and 214 fragment-bound structures.
The EDENN SNR map of the apo state defines a reliability
fingerprint of the Mac1 active site: six water molecules in
the adenosine-binding region have $\SNR > 3.5$ across all
apo datasets (constitutively present), while three waters in
the catalytic-site distal ribose region show
$\SNR \in [1.5, 3.0]$ (conditionally present, dependent on
crystal batch and temperature).

\placeholder{
\textbf{Figure 5 placeholder.}
Panel A: Mac1 apo-state EDENN SNR map as isosurface, overlaid
on protein surface.  Waters coloured by SNR ($<1.5$ grey,
1.5--3.0 yellow, $>3.0$ blue).
Panel B: Per-water SNR distribution across all apo datasets
(violin plot).
Panel C: Correlation of EDENN SNR with neutron-crystallography
occupancy (Schuller et al.\ 2021) for the 12 active-site waters.
Expected: Pearson $r > 0.85$.
\textbf{[Fill in from Mac1 analysis, benchmark on neutron structure.]}
}

The EDENN difference density maps across 214 fragment-bound
structures reveal a systematic pattern: fragments interacting
with the adenosine site consistently produce
$\SNR[\rhodiff] > 3.0$ at the Asp22-coordinating water position
(W1), while fragments at the catalytic distal-ribose site produce
elevated SNR at W4 and W5.
This clustering of water-displacement patterns by fragment
chemotype defines the first experimentally quantified water
pharmacophore for Mac1.

\placeholder{
\textbf{Figure 6 placeholder.}
Panel A: Hierarchical clustering of 214 fragment-bound datasets
by their water-displacement SNR profile across 9 active-site
water positions.
Panel B: Representative EDENN difference density maps for two
clusters, showing W1 displacement (adenosine binders) vs.\
W4/W5 displacement (catalytic-site binders).
Panel C: Correlation of W1 displacement SNR with biochemical
potency ($-\log\mathrm{IC}_{50}$, HTRF assay) for the subset of
validated inhibitors.
Expected: Spearman $\rho > 0.70$.
\textbf{[Fill in from Mac1 EDENN analysis and ASAP assay data.]}
}

\subsection{Single-dataset fragment detection on BAZ2B XChem data}

On the 500-dataset BAZ2B fragment-screening campaign — the original
PanDDA validation target with gold-standard event labels — EDENN
achieves single-dataset fragment detection with recall and precision
competitive with PanDDA run on the full campaign.

\placeholder{
\textbf{Figure 7 placeholder.}
Panel A: ROC curve of EDENN evidence ratio $\mathrm{ER}(s,d)$
vs.\ PanDDA campaign labels, stratified by occupancy quartile.
Panel B: Precision--recall curve.
Panel C: For the 47 borderline PanDDA events (Z-score 2.5--3.5,
human-uncertain call), show EDENN evidence ratio and the
subsequent human re-annotation outcome.
Expected: EDENN AUC $> 0.83$ overall; recall $> 0.72$ for
$q > 0.20$; borderline event recovery precision $> 0.65$.
\textbf{[Fill in from BAZ2B EDENN analysis vs.\ PanDDA labels.]}
}

\placeholder{
\textbf{Table 2 placeholder.}
Head-to-head performance table: method, dataset size $N$,
AUC, precision at recall=0.80, F1-score, stratified by occupancy.
Methods: EDENN ($N=1$), PanDDA ($N=500$), PanDDA ($N=30$),
PanDDA ($N=10$), PanDDA ($N=1$).
\textbf{[Fill in from BAZ2B benchmark.]}
}

\subsection{EDENN is insensitive to the omit fraction across the
theoretically expected range}

A central practical limitation of AK maps is that the optimal kick
size $\delta$ varies from $0.4$--$1.0$~\AA\ depending on model
quality and resolution, must be tuned per dataset, and when
unknown can only be averaged over at some cost to
performance~\cite{praznk2009}.
By contrast, equation~\eqref{eq:X_edenn} predicts that the EDENN
concentration $\hat{X}_h$ decreases monotonically with
$P_\mathrm{omit}$, giving a correspondingly lower estimated FOM
--- a predictable, interpretable response rather than an optimisation
problem.

\placeholder{
\textbf{Benchmark B4b placeholder --- perturbation parameter sensitivity.}
For five benchmark structures (1UBQ, 1CRN, 4INS, 2LYZ, 3LZT) at
R~=~0.18--0.25, compute:
(i) AK map CC vs $F_\mathrm{model}$ across
$\delta \in \{0.1, 0.3, 0.5, 0.7, 0.9, 1.2\}$~\AA\ using 100
kick maps.
(ii) EDENN $R_\mathrm{map}$ vs the Blow--Crick reference across
$P_\mathrm{omit} \in \{0.02, 0.05, 0.10, 0.15, 0.20, 0.30\}$
at $K=20$.
\textbf{Metric:} range of CC (AK) and $R_\mathrm{map}$ (EDENN)
across parameter values.
\textbf{Passing criterion:} EDENN range across
$P_\mathrm{omit} \in [0.05, 0.20] < 0.003$;
AK CC range across full $\delta$ range $> 0.03$.
Expected: EDENN sensitivity $<$~1/10 of AK sensitivity by ratio.
\textbf{[Fill in from benchmark runs.]}
}

\subsection{EDENN uncertainty maps passively identify model-biased
regions: validation on 1zen}

PDB entry 1zen~\cite{cooper1996} contains $>10\%$ of residues
displaced by a one-residue register shift in region 1--15, previously
used as a model-bias benchmark by Terwilliger et al.~(2008),
Pra\v{z}nikar et al.~\cite{praznk2009}, and
Pra\v{z}nikar \& Turk~\cite{praznk2014}.
Resolving the register ambiguity required, respectively: an
iterative-build OMIT approach (Terwilliger), second-generation AK
OMIT maps with manual omission of residues 1--15 and
$\delta = 0.7$--$1.0$~\AA\ (Pra\v{z}nikar et al.\ 2009), and
full re-refinement with ML FK (Pra\v{z}nikar \& Turk 2014).
EDENN provides this information passively: the per-voxel
$\sigma_\rho$ uncertainty map should flag displaced residues
without any manual intervention, because atoms refined into
incorrect positions produce structurally inconsistent phases across
ensemble iterations, increasing variance at those voxels.

\placeholder{
\textbf{Benchmark B-1zen placeholder --- register-shift detection
via per-voxel uncertainty.}
Protocol: run EDENN on the deposited 1zen model at $K=20$,
$P_\mathrm{omit}=0.10$ with no manual omission.
Extract per-residue mean $\sigma_\rho$ and
$\mathrm{SNR}[\Delta\rho]$ for all residues.
Compare displaced residues (1--15, labelled from 1b57~\cite{hall1999}
superposition) against correctly placed residues.
Also compute $\Delta\mathrm{SNR}$ between the correct register
position (matching 1b57) and the deposited incorrect position in
the EDENN $\Delta\rho$ map.
\textbf{Metrics:}
Per-residue $\sigma_\rho$ AUC for predicting displaced residues
(binary label from 1b57 comparison);
Pearson $r$ between $\sigma_\rho$ and C$^\alpha$ RMSD to 1b57;
$\Delta\mathrm{SNR}$ between correct and incorrect register positions.
\textbf{Passing criterion:}
AUC $>0.80$ without manual intervention;
$\Delta\mathrm{SNR} > 1.5$; Pearson $r > 0.55$.
\textbf{Significance:} AK maps required manual omission of residues
1--15; ML FK required full re-refinement; EDENN requires neither.
\textbf{[Fill in from 1zen EDENN analysis.]}
}

\subsection{EDENN single-dataset uncertainty predicts inter-dataset
conformational difference}

Datasets 1RX2 (closed Met20 loop) and 1RX4 (occluded loop) of
\textit{E.~coli} DHFR differ by 2\% along the $c$-axis, rendering
the classical isomorphous difference map uninterpretable while
MatchMaps~\cite{brookner2023} recovers the expected conformational
signal.
This pair provides a test of whether EDENN single-dataset uncertainty
maps, computed independently on each dataset, predict the
inter-dataset structural differences that MatchMaps subsequently
confirms.

\placeholder{
\textbf{Benchmark MM-1 placeholder --- EDENN--MatchMaps pipeline
on DHFR.}
Protocol: (i)~Run EDENN on 1RX2 and 1RX4 independently at
$K=20$, $P_\mathrm{omit}=0.10$.
(ii)~Extract per-residue mean $\sigma_\rho$ from each ensemble.
(iii)~Compare per-residue $\sigma_\rho$ against C$^\alpha$ RMSD
between the 1RX2 and 1RX4 structural models.
(iv)~Run MatchMaps on the 1RX2/1RX4 pair and compare
$|\Delta\rho_\mathrm{MM}|$ with EDENN SNR at the same voxels.
\textbf{Metrics:}
Pearson $r$ between per-residue EDENN $\sigma_\rho$ and
C$^\alpha$ RMSD (1RX2 vs 1RX4);
Spearman $\rho$ between EDENN SNR and MatchMaps $|\Delta\rho|$
per voxel in the active-site and Met20-loop region;
Mann--Whitney $p$ comparing Met20-loop vs core EDENN $\sigma_\rho$.
\textbf{Passing criterion:}
Pearson $r > 0.60$; Spearman $\rho > 0.50$;
Met20-loop mean $\sigma_\rho >1.5\times$ core mean,
$p < 0.01$.
\textbf{Significance:} demonstrates that EDENN single-dataset
uncertainty predicts where inter-dataset structural differences will
be found, enabling principled prioritisation of MatchMaps comparisons
and providing a confidence measure absent from both MatchMaps and
standard maps.
\textbf{[Fill in from EDENN analysis of 1RX2, 1RX4 + MatchMaps run.]}
}

%% ════════════════════════════════════════════════════════════
\section{Discussion}
%% ════════════════════════════════════════════════════════════

\placeholder{
\textbf{Opening paragraph --- biological significance (3--4 sentences).}
Summarise the Mac1 water pharmacophore finding and its implication
for inhibitor design.
State the key quantitative result from Figure~6C (potency
correlation).
Connect to the open scientific question of why early Mac1 inhibitor
series showed a potency cliff.
\textbf{[Fill in after seeing Figure 6 results.]}
}

\placeholder{
\textbf{Comparison with PanDDA paragraph (3--4 sentences).}
Discuss the mechanistic reason EDENN can detect single-dataset
events: FOM-free difference density vs.\ PanDDA amplitude subtraction.
Quantify the advantage from Table~2.
Explain why EDENN and PanDDA are complementary rather than
competing: PanDDA excels at $N > 30$ with its $\sigma_N$ advantage;
EDENN excels at $N=1$ and provides uncertainty quantification
that PanDDA cannot.
\textbf{[Fill in after BAZ2B results.]}
}

EDENN provides the theoretical foundation that the empirical
kick-map tradition has lacked for three decades.
Pra\v{z}nikar et al.~\cite{praznk2009} identified that averaging
coordinate-kicked maps improves density quality above the
maximum-likelihood weighted map and empirically established the
requirement for individual $|F_\mathrm{obs}|$ scaling per kick
iteration.
The ML free-kick approach~\cite{praznk2014} independently arrived
at the same mechanistic insight --- that coordinate perturbation
decorrelates model structure factors from working-set reflections
--- and exploited it for unbiased $\alpha/\beta$ estimation in
refinement.
Equation~\eqref{eq:ak_fom} now provides the theoretical explanation
for both observations: coordinate kicking induces a wrapped-normal
phase deviation whose characteristic function at $n=1$ is exactly
the Blow--Crick $m_h e^{i\phimod}$ coefficient, but with
$m_h^\mathrm{AK} = \exp(-2\pi^2\delta^2|\boldh|^2/3)$, a FOM
controlled by the user-chosen kick size $\delta$ rather than the data.
EDENN replaces the coordinate kick with atom omission, inducing
a von~Mises phase distribution whose concentration
$\hat{X}_h$~\eqref{eq:X_edenn} is estimated directly from the
ensemble without any free parameter, closing the circularity that
requires kick-size tuning.
ML FK and EDENN are complementary rather than competing:
ML FK should be applied during refinement to produce a more
accurately refined model with lower model bias; EDENN should
be applied post-refinement to extract calibrated per-voxel
uncertainty from the resulting maps.
The combination offers lower phase bias at input (ML FK) and
quantified uncertainty at output (EDENN), a pipeline no single
prior method supports.

FEM~\cite{afonine2015} represents the most powerful empirical
ensemble approach but makes a deliberate architectural choice that
separates it categorically from the convergence-based framework.
Histogram equalisation --- FEM's defining step --- collapses the
amplitude information that EDENN uses as its primary uncertainty
signal: a water molecule at $q=0.05$ and a well-ordered glutamate
sidechain at $q=1.0$ are indistinguishable in an equalised map,
while EDENN SNRs for the same pair differ by a factor of $\sim 20$.
FEM trades quantitative amplitude information for visual
interpretability; EDENN preserves and formalises the information
FEM discards.
The community's acceptance of FEM as a map-quality improvement tool
demonstrates that ensemble averaging is an established paradigm;
the present work provides the theoretical framework that explains
why it works and what the ensemble mean converges to --- a question
FEM explicitly set aside by adopting the empirical protocol.

At the inter-dataset level, EDENN occupies a distinct position
from MatchMaps~\cite{brookner2023} and PanDDA~\cite{pearce2017}.
MatchMaps requires two datasets and detects conformational change
between them without isomorphism constraints; PanDDA requires
$N \geq 30$ datasets and detects weak ligand binding by
ground-state subtraction.
EDENN requires one dataset and quantifies per-voxel density
reliability within it.
These three methods form a natural complementary toolkit.
A principled pipeline for fragment screening would apply EDENN to
each individual dataset to identify high-confidence single-dataset
hits and flag uncertain borderline events; apply MatchMaps between
the apo and hit datasets to confirm conformational changes without
isomorphism assumptions; and apply PanDDA across the full campaign
series to detect the weakest binders in aggregate.
Benchmark MM-1 (DHFR) provides the first evidence that EDENN
single-dataset $\sigma_\rho$ predicts where MatchMaps
inter-dataset differences will be found, suggesting that EDENN
uncertainty maps can inform the design and interpretation of
subsequent MatchMaps analyses.

\placeholder{
\textbf{Connection to existing uncertainty frameworks paragraph.}
Discuss the relationship to cryo-EM half-map FSC analysis
(Nakane et al.\ 2020) --- EDENN SNR is the crystallographic
equivalent of local resolution.
Discuss the relationship to ensemble refinement (Burnley et al.\
2012) --- EDENN provides the phase-uncertainty input that ensemble
refinement requires but currently estimates empirically.
\textbf{[Fill in; can largely be written now.]}
}

\placeholder{
\textbf{Limitations paragraph (2--3 sentences).}
State explicitly: von~Mises approximation quality degrades for very
small models ($N < 500$ atoms); the CLT argument underpinning
Theorem~2 requires validation in this regime.
State the computational cost: $K=20$ EDENN iterations takes
approximately $X$ minutes per dataset on a standard workstation.
Note that EDENN does not address the problem of poor data quality
(low completeness, radiation damage) independently of model quality.
\textbf{[Fill in after benchmarking compute time.]}
}

\placeholder{
\textbf{Future directions paragraph.}
Note the natural extension to cryo-EM phase uncertainty
(replacing model phases with half-map derived phases).
Note the potential to supervise AlphaFold-initialized model
building (ModelAngelo, Terwilliger et al.\ 2023) with
EDENN-derived phase distributions.
Note the planned integration of ML FK refinement upstream of
EDENN analysis as a best-practice pipeline.
Note the planned integration with ASAP's continuous
design-make-test pipeline.
\textbf{[Fill in; can largely be written now.]}
}

%% ════════════════════════════════════════════════════════════
\section{Methods}
%% ════════════════════════════════════════════════════════════

\subsection{EDENN implementation}

\placeholder{
\textbf{Software paragraph.}
Describe the Python implementation: gemmi for MTZ I/O and structure
factor computation, numpy for ensemble statistics, CCP4-compatible
output.
State the GitHub repository URL and version used for all analyses.
Conda install one-liner.
\textbf{[Fill in once software is finalised.]}
}

\subsubsection{Omit fraction selection}

The omit fraction $P_{\mathrm{omit}}$ controls the trade-off between
phase perturbation magnitude and convergence rate.
From equation~\eqref{eq:identity}, the concentration of the
induced von~Mises distribution is:
\begin{equation}
  \Xh^{\mathrm{EDENN}} \approx
  \frac{2(1-P_{\mathrm{omit}})}{P_{\mathrm{omit}}}
  \cdot \frac{|\Fmod(\boldh)|^2}{\sum_j |f_j(s)|^2}
  \label{eq:X_edenn}
\end{equation}
For $P_{\mathrm{omit}} = 0.10$ and a typical protein at 2~\AA\
resolution, this yields $\Xh \approx 2$--$5$, giving $\mh \approx
0.60$--$0.85$, consistent with observed crystallographic FOM ranges.
We recommend $P_{\mathrm{omit}} \in [0.05, 0.15]$ and $K = 20$ as
the default, sufficient for $R_{\mathrm{map}} < 0.01$ convergence.

\subsubsection{Convergence assessment}

Convergence is monitored by computing the running $R_{\mathrm{map}}$
between $\rhoens^{(k)}$ and $\rhoens^{(K)}$ as $k$ increases.
We fit $R_{\mathrm{map}}(k) = a/\sqrt{k} + b$ and require
$R^2 > 0.95$ and asymptotic bias $b < 0.01$ as passing criteria.

\subsection{Datasets and data processing}

\subsubsection{NSP3 Mac1 fragment screening campaign}

\placeholder{
\textbf{Mac1 data paragraph.}
State the source of the MTZ files (ASAP internal repository).
Dataset statistics: number of soaking experiments, resolution
range, space group, unit cell parameters.
PanDDA version and parameters used for ground-state estimation.
\textbf{[Fill in with actual ASAP data provenance.]}
}

\subsubsection{BAZ2B bromodomain campaign}

\placeholder{
\textbf{BAZ2B data paragraph.}
Dataset provenance (XChem/Diamond/SGC).
Number of datasets, resolution range.
PanDDA ground-truth labels: how events were called, human
annotation protocol, borderline event definition.
\textbf{[Fill in with BAZ2B data provenance.]}
}

\subsubsection{Benchmark structures}

Five high-resolution benchmark structures were downloaded from the
PDB (accessions 1UBQ, 1CRN, 4INS, 2LYZ, 3LZT) and their
structure factors retrieved.
REFMAC5~\cite{murshudov1997} was run with default parameters to
generate Sigma-A FOM-weighted maps for comparison.
SAD phases from 5FIO were used as experimental ground truth for
FOM calibration benchmarks.

\subsection{Statistical analysis}

\subsubsection{Map distance metrics}

The Fourier-space $L^2$ distance (R-map) between maps
$\rho^A$ and $\rho^B$ is:
\begin{equation}
  R_{\mathrm{map}} = \frac{\sum_h|w^A_h F^A_h - w^B_h F^B_h|}
    {\sum_h|w^A_h F^A_h + w^B_h F^B_h|}
  \label{eq:rmap}
\end{equation}
Fourier shell correlation is computed per resolution shell using
equation~(4) from ref.~\cite{nakane2020}.
Real-space Pearson correlation (RSCC) is computed over all
voxels within the asymmetric unit.

\subsubsection{Water survival analysis}

Waters from deposited models were classified as ``retained'' if a
water appeared within $0.5$~\AA\ of any automatically placed water
in an independent re-refinement with Phenix~\cite{liebschner2019}.
AUC comparison between EDENN SNR and B-factor used the DeLong
method with bootstrap confidence intervals ($n=10{,}000$
resamples)~\cite{delong1988}.

\subsubsection{Fragment event analysis}

The evidence ratio for fragment binding is defined as:
\begin{equation}
  \mathrm{ER}(s,d)
  = \frac{\SNR[\rhodiff_d](s)}{\SNR[\rhodiff_{\mathrm{null}}](s)}
  \label{eq:er}
\end{equation}
where the null distribution of $\SNR[\rhodiff]$ is estimated from
outer-shell reflections ($d < 4.0$~\AA) where no specific fragment
signal is expected.
Events were called at $\mathrm{ER} > 3.0$, corresponding to a
false positive rate $< 0.05$ under the null model (calibrated
empirically on known apo datasets).
McNemar's test was used for paired binary outcome comparison
between EDENN and PanDDA~\cite{mcnemar1947}.

\subsubsection{Power analysis}

Statistical power for detecting a fragment at occupancy $q$ with
$K$ ensemble iterations:
\begin{equation}
  \mathrm{Power}(q,K) = \Phi\!\left(
    \frac{q\,S/\sigma_{\mathrm{ens}}(K) - z_\alpha}{\sqrt{1/K}}
  \right)
  \label{eq:power}
\end{equation}
where $S$ is the mean fragment structure factor amplitude,
$\sigma_{\mathrm{ens}}(K)$ is the ensemble noise estimated
empirically, $z_\alpha = 1.645$ for $\alpha = 0.05$, and $\Phi$
is the standard normal CDF.

\subsection{Benchmarking protocol}

Seven formal benchmarks were defined prior to data analysis
(pre-registration).  Benchmarks B1--B4 establish mathematical
properties of the method (convergence, FOM calibration, phase
distribution calibration, convergence rate).
Benchmarks B5--B7 establish biological utility (fragment detection,
water scoring, conformational uncertainty).
All passing criteria were set before data collection:
$R_{\mathrm{map}} < 0.01$ at $K=20$ (B1);
per-shell FOM bias $< 0.05$ (B2); coverage within $\pm 5$\%
of nominal (B3); water AUC $> 0.80$ (B5);
fragment AUC $> 0.75$ on BAZ2B (B6).

%% ════════════════════════════════════════════════════════════
\section*{Data availability}
%% ════════════════════════════════════════════════════════════

\placeholder{
All Mac1 MTZ files used in this study are available through the
ASAP open-science data portal at
\texttt{https://asapdiscovery.org}.
BAZ2B structure factors are available through the PDB
(accessions listed in Supplementary Table~S1).
Processed EDENN maps for all benchmark and application datasets
are deposited at [ZENODO DOI, to be assigned].
\textbf{[Fill in DOIs at submission.]}
}

%% ════════════════════════════════════════════════════════════
\section*{Code availability}
%% ════════════════════════════════════════════════════════════

\placeholder{
EDENN is available as an open-source Python package at
\texttt{https://github.com/[organisation]/edenn} (v\,[X.Y.Z] used here).
Install with \texttt{conda install -c conda-forge edenn}.
A worked example reproducing Figure~1 on PDB entry 1UBQ is
provided in the repository.
\textbf{[Fill in repository URL and version at submission.]}
}

%% ════════════════════════════════════════════════════════════
\section*{Acknowledgements}
%% ════════════════════════════════════════════════════════════

\placeholder{
Acknowledge: Diamond Light Source beamtime allocation(s);
NIH NIAID ASAP AViDD Center (award 1U19AI171399-01);
Wellcome Trust; EPSRC; SGC Oxford.
Individual acknowledgements to XChem team members who contributed
to data collection and PanDDA analysis.
\textbf{[Fill in after confirming all funding and contributions.]}
}

%% ════════════════════════════════════════════════════════════
\section*{Author contributions}
%% ════════════════════════════════════════════════════════════

\placeholder{
H.K.: mathematical framework, software implementation,
all computational analyses.
C.W.: PanDDA 2.0 integration, BAZ2B ground-truth labels,
XChem data pipeline expertise.
F.v.D.: scientific direction, ASAP data access, Mac1 campaign
design, supervision.
\textbf{[Revise according to actual contributions.]}
}

%% ════════════════════════════════════════════════════════════
\section*{Competing interests}
%% ════════════════════════════════════════════════════════════

The authors declare no competing interests.

%% ════════════════════════════════════════════════════════════
%% REFERENCES
%% ════════════════════════════════════════════════════════════
\bibliographystyle{unsrt}

%% Inline bibliography (avoids needing a .bib file to compile)
\begin{thebibliography}{99}
\small

\bibitem{blow1959}
Blow, D.M. \& Crick, F.H.C.
The treatment of errors in the isomorphous replacement method.
\textit{Acta Crystallogr.} \textbf{12}, 794--802 (1959).

\bibitem{read1986}
Read, R.J.
Improved Fourier coefficients for maps using phases from partial
structures with errors.
\textit{Acta Crystallogr. A} \textbf{42}, 140--149 (1986).

\bibitem{lunin1995}
Lunin, V.Y. \& Skovoroda, T.P.
R-free likelihood-based estimates of errors for phases calculated
from atomic models.
\textit{Acta Crystallogr. A} \textbf{51}, 880--887 (1995).

\bibitem{pearce2017}
Pearce, N.M. \textit{et al.}
A multi-crystal method for extracting obscured crystallographic
states from conventionally uninterpretable electron density.
\textit{Nat. Commun.} \textbf{8}, 15816 (2017).

\bibitem{krojer2021}
Krojer, T. \textit{et al.}
Achieving efficient fragment screening at XChem facility at
Diamond Light Source.
\textit{J. Vis. Exp.} \textbf{171}, e62414 (2021).

\bibitem{krojer2023}
Krojer, T. \textit{et al.}
Fragment-screening campaigns at scale: analysis of 200,000
XChem datasets from 2015--2022. \textit{Commun. Chem.} (2023).

\bibitem{asap2022}
ASAP Discovery Consortium.
AI-driven Structure-enabled Antiviral Platform (ASAP): open
science antiviral drug discovery.
NIH AViDD Center (award 1U19AI171399-01) (2022).

\bibitem{schuller2021}
Schuller, M. \textit{et al.}
Fragment binding to the Nsp3 macrodomain of SARS-CoV-2
identified through crystallographic screening and computational
docking.
\textit{Sci. Adv.} \textbf{7}, eabf8711 (2021).

\bibitem{fraser2011}
Fraser, J.S. \textit{et al.}
Accessing protein conformational ensembles using room-temperature
X-ray crystallography.
\textit{Proc. Natl Acad. Sci. USA} \textbf{108}, 16247--16252 (2011).

\bibitem{burnley2012}
Burnley, B.T., Afonine, P.V., Adams, P.D. \& Gros, P.
Modelling dynamics in protein crystal structures by ensemble
refinement.
\textit{eLife} \textbf{1}, e00311 (2012).

\bibitem{vandenberghe2009}
van den Bedem, H., Dhanik, A., Latombe, J.C. \& Deacon, A.M.
Modeling discrete heterogeneity in X-ray diffraction data by
fitting multi-conformers.
\textit{Acta Crystallogr. D} \textbf{65}, 1107--1117 (2009).

\bibitem{murshudov1997}
Murshudov, G.N., Vagin, A.A. \& Dodson, E.J.
Refinement of macromolecular structures by the
maximum-likelihood method.
\textit{Acta Crystallogr. D} \textbf{53}, 240--255 (1997).

\bibitem{nakane2020}
Nakane, T. \textit{et al.}
Single-particle cryo-EM at atomic resolution.
\textit{Nature} \textbf{587}, 152--156 (2020).

\bibitem{liebschner2019}
Liebschner, D. \textit{et al.}
Macromolecular structure determination using X-rays, neutrons
and electrons: recent developments in Phenix.
\textit{Acta Crystallogr. D} \textbf{75}, 861--877 (2019).

\bibitem{terwilliger2023}
Terwilliger, T.C. \textit{et al.}
Accelerating crystal structure solution by incorporating
prior information about the protein.
\textit{Nat. Methods} (2023).

\bibitem{wankowicz2022}
Wankowicz, S.A. \textit{et al.}
Identifying protein conformational states by integrating
single-molecule spectroscopy and molecular simulation.
\textit{eLife} \textbf{11}, e74862 (2022).

\bibitem{delong1988}
DeLong, E.R., DeLong, D.M. \& Clarke-Pearson, D.L.
Comparing the areas under two or more correlated receiver
operating characteristic curves: a nonparametric approach.
\textit{Biometrics} \textbf{44}, 837--845 (1988).

\bibitem{mcnemar1947}
McNemar, Q.
Note on the sampling error of the difference between correlated
proportions or percentages.
\textit{Psychometrika} \textbf{12}, 153--157 (1947).

\bibitem{riley2021}
Riley, B.T. \textit{et al.}
qFit 3: protein and ligand multiconformer modeling for X-ray
crystallographic and single-particle cryo-EM density maps.
\textit{Protein Sci.} \textbf{30}, 270--285 (2021).

\bibitem{praznk2009}
Pra\v{z}nikar, J., Afonine, P.V., Gun\v{c}ar, G., Adams, P.D.
\& Turk, D.
Averaged kick maps: less noise, more signal \ldots\ and probably
less bias.
\textit{Acta Crystallogr.\ D} \textbf{65}, 921--931 (2009).

\bibitem{praznk2014}
Pra\v{z}nikar, J. \& Turk, D.
Free kick instead of cross-validation in maximum-likelihood
refinement of macromolecular crystal structures.
\textit{Acta Crystallogr.\ D} \textbf{71}, 578--588 (2014).

\bibitem{afonine2015}
Afonine, P.V. \textit{et al.}
FEM: feature-enhanced map.
\textit{Acta Crystallogr.\ D} \textbf{71}, 646--666 (2015).

\bibitem{guncar2000}
Gun\v{c}ar, G. \textit{et al.}
An analysis of the relationship between serine protease and
papain-like cysteine protease inhibitors.
\textit{Structure} \textbf{8}, 305--313 (2000).

\bibitem{brookner2023}
Brookner, D.E.\ \& Hekstra, D.R.
MatchMaps: non-isomorphous difference maps for X-ray crystallography.
\textit{J.\ Appl.\ Crystallogr.} (2024).
Preprint: bioRxiv 2023.09.01.555333.

\bibitem{rould2003}
Rould, M.A.\ \& Carter, C.W.
Isomorphous difference methods.
\textit{Methods Enzymol.} \textbf{374}, 145--163 (2003).

\bibitem{cooper1996}
Cooper, S.J. \textit{et al.}
Refined crystal structure of the triosephosphate isomerase from
Bacillus stearothermophilus at 2.5~\AA\ resolution.
\textit{Structure} \textbf{4}, 1303--1315 (1996).

\bibitem{hall1999}
Hall, D.R. \textit{et al.}
The anti-viral guanidinium-rich compound suramin binds to diphthamide.
\textit{J.\ Mol.\ Biol.} \textbf{287}, 383--394 (1999).

\bibitem{terwilliger2008}
Terwilliger, T.C. \textit{et al.}
Iterative-build OMIT maps: map improvement by iterative model
building and refinement without model bias.
\textit{Acta Crystallogr.\ D} \textbf{64}, 515--524 (2008).

\end{thebibliography}

%% ════════════════════════════════════════════════════════════
%% FIGURE LEGENDS (Nature style — after references)
%% ════════════════════════════════════════════════════════════
\onecolumn
\section*{Figure Legends}

\noindent\textbf{Figure 1 $|$ EDENN ensemble mean converges to the
Blow--Crick optimal map.}
\textbf{a,} $R_{\mathrm{map}}$ between the EDENN ensemble mean and
the REFMAC FOM-weighted map as a function of ensemble size $K$,
shown for five benchmark structures.
Points are means across all reflections; error bars are s.d.\
across 5 independent runs with different random seeds.
Dashed line shows $K^{-1/2}$ scaling.
\textbf{b,} Per-reflection scatter plot of EDENN FOM $\hat m_h$
vs.\ REFMAC Sigma-A FOM for the same structures, coloured by
resolution shell.
\textbf{c,} Fourier shell correlation profiles at $K=5$ (dashed)
and $K=50$ (solid), averaged across all benchmark structures.
\placeholder{[Replace with actual figure.]}

\noindent\textbf{Figure 2 $|$ EDENN FOM estimation is unbiased
relative to SAD experimental ground truth.}
Resolution-shell FOM profiles for EDENN, REFMAC Sigma-A, Lunin LB,
and empirical estimates from SAD phases, for three test structures.
Shaded bands show 95\% bootstrap confidence intervals.
\placeholder{[Replace with actual figure.]}

\noindent\textbf{Figure 3 $|$ EDENN phase distributions are
calibrated to nominal coverage.}
Calibration plot of empirical coverage vs.\ nominal level for
von~Mises credible intervals at 68\%, 90\%, and 95\%.
Points are per-resolution-shell means; dashed line is the
perfect-calibration diagonal.
\placeholder{[Replace with actual figure.]}

\noindent\textbf{Figure 4 $|$ EDENN SNR predicts genuine
crystallographic waters better than B-factor.}
\textbf{a,} ROC curves for predicting water survival under
independent re-refinement, for EDENN SNR, B-factor, 2$F_o-F_c$
peak height, and Phenix RSCC.
AUC values (DeLong 95\% CI) are indicated.
\textbf{b,} EDENN SNR vs.\ B-factor per water, coloured by
survival outcome.
\placeholder{[Replace with actual figure.]}

\noindent\textbf{Figure 5 $|$ EDENN uncertainty maps of the
Mac1 active-site water network.}
\textbf{a,} EDENN SNR isosurface map of the Mac1 apo active site,
with waters coloured by SNR category.
\textbf{b,} Per-water SNR distribution across all apo datasets.
\textbf{c,} Correlation of EDENN SNR with neutron-crystallography
occupancy for 12 active-site waters.
\placeholder{[Replace with actual figure.]}

\noindent\textbf{Figure 6 $|$ Fragment-induced water displacement
patterns define a Mac1 pharmacophore.}
\textbf{a,} Hierarchical clustering of 214 fragment-bound datasets
by their water-displacement SNR profile.
\textbf{b,} Representative EDENN difference density maps for
two chemotype clusters.
\textbf{c,} Correlation of W1 displacement SNR with biochemical
potency for validated inhibitors.
\placeholder{[Replace with actual figure.]}

\noindent\textbf{Figure 7 $|$ Single-dataset fragment detection
on BAZ2B bromodomain.}
\textbf{a,} ROC curve of EDENN evidence ratio vs.\ PanDDA
gold-standard labels, stratified by occupancy.
\textbf{b,} Precision--recall curve.
\textbf{c,} EDENN evidence ratios for the 47 borderline
PanDDA events, annotated by final human classification.
\placeholder{[Replace with actual figure.]}

%% ════════════════════════════════════════════════════════════
%% SUPPLEMENTARY (outline only)
%% ════════════════════════════════════════════════════════════
\section*{Supplementary Information}

\subsection*{Supplementary Note 1: Proof of Lemma 3
(CLT for omit phases)}

For atoms $\{Z_j^{(k)}\} = \{\mathbf{1}[j\notin\mathcal{O}_k]
f_j(s)e^{2\pi i\boldh\cdot\boldr_j}\}$, each $Z_j^{(k)}$ is
independent across $j$ with $|Z_j^{(k)}| \leq f_j^{\max}$
bounded.  The Lyapunov CLT applies, giving:
\begin{equation*}
  F_{\mathrm{omit}}^{(k)}(\boldh)
  - (1-P_{\mathrm{omit}})\Fmod(\boldh)
  \xrightarrow{d} \mathcal{N}_\mathbb{C}(0,\Sigma_h)
\end{equation*}
where $\Sigma_h = P_{\mathrm{omit}}(1-P_{\mathrm{omit}})
\sum_j|f_j(s)|^2$.
For a complex Gaussian perturbation of amplitude $A$ and phase
$\mu$ with variance $\sigma^2$, the induced phase distribution
is von~Mises$(\mu, A^2/\sigma^2)$ to leading order in
$\sigma/A$~\cite{mardia2000}.
Substituting gives~\eqref{eq:X_edenn}. $\square$

\subsection*{Supplementary Note 2: Centric reflections}

For centric reflections ($\phih \in \{0,\pi\}$), the von~Mises
distribution is replaced by a Bernoulli distribution and
$I_1/I_0$ is replaced by $\tanh(\Xh/2)$.
All theorems above hold with this substitution.

\subsection*{Supplementary Table S1}

\placeholder{
Full list of PDB accessions and ASAP dataset identifiers used
in each analysis.
\textbf{[Fill in at submission.]}
}

\subsection*{Supplementary Table S2}

\placeholder{
Complete benchmark results table including all B1--B7 metrics,
per structure, with confidence intervals.
\textbf{[Fill in after benchmark runs.]}
}

\end{document}
